% TT07XX1C
% Polytrauma, allgemein
% 14.0
% 2013-11-30
:ref:`m-polytrauma`
EE{Taktik}: \Ca{Vitale Bedrohung! Zeitkritisch!}
Rasche Bergung, Stabilisierung
der Vitalfunktionen, Immobilisierung, rasche Versorgung von
*gefährlichen* Verletzungen, rascher Transport}


\begin{itemize}
  -   |TxMassVitMKBes|
  \begin{itemize}
    -   Lagerung: situationsgerecht, ggfs.
    Wirbelsäulenimmobilisation, bzw. entsprechend den
    Verletzungen und der Herz-Kreislauf-Situation
%     -- Prioritäten absteigend:
%     {
%     \scriptsize
%     \begin{enumerate}
%       -   Atem- und Kreislaufstillstand $\rightarrow$ Rückenlage
%       -   Bewusstlosigkeit, nicht intubiert $\rightarrow$ stabile Seitenlage 
%       (bei Thorax- und/oder Extremitätentrauma auf verletzte Seite legen)
%       -   Atemnot bei Thoraxtrauma $\rightarrow$ Oberkörper hoch (45\textdegree)
%       -   Abdominaltrauma $\rightarrow$ Schonhaltung, Knierolle
%       -   \Abk{SHT} $\rightarrow$ Oberkörper hoch (30\textdegree)
%       -   Wirbelsäulentrauma $\rightarrow$ Flachlagerung auf Vakuummatratze
%       -   Extremitätentrauma $\rightarrow$ Extremität schienen und fixieren
%       -   Schock $\rightarrow$ Hochlagerung der Beine
%       durch schräg stellen des Tragentisches
%     \end{enumerate}
%     }
  \end{itemize}
  -   Gefährliche Verletzungen versorgen (starke
  Blutungen, Beckengurt, \ldots) 
  -   **Wärmeerhalt**: Der Wärmeerhalt hat massiven Einfluss auf das
  Überleben!
%   -   Wenn einigermaßen stabil, gründlicher  Traumacheck
  -   Verletzungsmuster/Unfallmechanismus erheben $\rightarrow$ zu 
  erwartende Komplikationen einschätzen
  -   \Z{Notarzt: Schmerztherapie, Narkose, Intubation und
  Beatmung}
  -   Transportentscheidung:
  \begin{itemize}
    -   Transportentscheidung: Schockraum, mit
    Voranmeldung (Aviso) über Leitstelle
    -   Hubschrauber-Transport erwägen
  \end{itemize}
\end{itemize}